% Options for packages loaded elsewhere
\PassOptionsToPackage{unicode}{hyperref}
\PassOptionsToPackage{hyphens}{url}
\PassOptionsToPackage{space}{xeCJK}
%
\documentclass[
]{ctexart}
\usepackage{amsmath,amssymb}
\usepackage{iftex}
\ifPDFTeX
  \usepackage[T1]{fontenc}
  \usepackage[utf8]{inputenc}
  \usepackage{textcomp} % provide euro and other symbols
\else % if luatex or xetex
  \usepackage{unicode-math} % this also loads fontspec
  \defaultfontfeatures{Scale=MatchLowercase}
  \defaultfontfeatures[\rmfamily]{Ligatures=TeX,Scale=1}
\fi
\usepackage{lmodern}
\ifPDFTeX\else
  % xetex/luatex font selection
  \ifXeTeX
    \usepackage{xeCJK}
    \setCJKmainfont[]{Songti SC}
          \fi
  \ifLuaTeX
    \usepackage[]{luatexja-fontspec}
    \setmainjfont[]{Songti SC}
  \fi
\fi
% Use upquote if available, for straight quotes in verbatim environments
\IfFileExists{upquote.sty}{\usepackage{upquote}}{}
\IfFileExists{microtype.sty}{% use microtype if available
  \usepackage[]{microtype}
  \UseMicrotypeSet[protrusion]{basicmath} % disable protrusion for tt fonts
}{}
\makeatletter
\@ifundefined{KOMAClassName}{% if non-KOMA class
  \IfFileExists{parskip.sty}{%
    \usepackage{parskip}
  }{% else
    \setlength{\parindent}{0pt}
    \setlength{\parskip}{6pt plus 2pt minus 1pt}}
}{% if KOMA class
  \KOMAoptions{parskip=half}}
\makeatother
\usepackage{xcolor}
\usepackage[margin=2.2cm]{geometry}
\setlength{\emergencystretch}{3em} % prevent overfull lines
\providecommand{\tightlist}{%
  \setlength{\itemsep}{0pt}\setlength{\parskip}{0pt}}
\setcounter{secnumdepth}{5}
\ifLuaTeX
  \usepackage{selnolig}  % disable illegal ligatures
\fi
\IfFileExists{bookmark.sty}{\usepackage{bookmark}}{\usepackage{hyperref}}
\IfFileExists{xurl.sty}{\usepackage{xurl}}{} % add URL line breaks if available
\urlstyle{same}
\hypersetup{
  pdftitle={Multi-Linear Regression},
  pdfauthor={Rui Zhou},
  hidelinks,
  pdfcreator={LaTeX via pandoc}}

\title{Multi-Linear Regression}
\author{Rui Zhou}
\date{}

\begin{document}
\maketitle

\renewcommand*\contentsname{目录}
{
\setcounter{tocdepth}{3}
\tableofcontents
}
\texttt{\{r\ warning=FALSE\}\ library(dplyr)}

The \texttt{df} contains information about salary, beginning salary,
education year, job category, previous working experience and so on. The
relationship between salary and beginning salary, education, to name a
few, is of interest.

\texttt{\{r\}\ df\ =\ readxl::read\_excel("data02-01.xlsx",\ sheet\ =\ 1)}

\hypertarget{pre-inspection}{%
\section{Pre-Inspection}\label{pre-inspection}}

\hypertarget{linear-relationship}{%
\subsection{linear relationship}\label{linear-relationship}}

Use \texttt{cor()} to carry the correlation test, the default method of
which is
\texttt{method\ =\ \textquotesingle{}pearson\textquotesingle{}}. The
result will be presented as a correlation matrix.

However, obviously the correlation matrix won't give you information
about p-value (or say, significance).

Good thing is that the check of linear relationship is the first and
rough step for further regression, maybe other more information (about
p-value or CIs) is unnecessary. Another good thing is that the
\texttt{cor()} is from base R, no extra requirement to library packages.

\texttt{\{r\ comment=\textquotesingle{}\textquotesingle{}\}\ df\ \textbar{}\textgreater{}\ select(salary,educ,salbegin,jobtime,prevexp)\ \textbar{}\textgreater{}\ cor()}

Well, it's clumsy to compress the target data (variables) into a matrix
to take the test and get the correlation matrix. Just fine to know a bit
about that.
\texttt{\{r\ comment=\textquotesingle{}\textquotesingle{}\}\ \#\ first\ create\ an\ empty\ matrix,\ primarily\ filled\ with\ NA\ m\ =\ matrix(nrow=length(salary),ncol=5)\ m{[},1{]}=salary\ m{[},2{]}=educ\ m{[},3{]}\ =\ salbegin\ m{[},4{]}=jobtime\ \ \ m{[},5{]}=prevexp\ \#check\ the\ correlation\ matrix\ cor(m)}

Additionally, I'd like to use \texttt{bruceR::Corr()} for the
correlation matrix, which is visible for both in amount and in
significance. Moreover, it's fancy!

For simplicity, the two decimals by default is well enough. If more
accuracy is needed, you can set the parameter \texttt{digits\ =}.

\texttt{\{r\ comment=\textquotesingle{}\textquotesingle{}\}\ df\ \textbar{}\textgreater{}\ select(salary,educ,salbegin,jobtime,prevexp)\ \textbar{}\textgreater{}\ bruceR::Corr(digits\ =\ 2)}

\hypertarget{multicollinearity}{%
\subsection{multicollinearity}\label{multicollinearity}}

Use \texttt{car::vif(model)} and get the result directly.

For the convenience of narration, check of multicollinearity is
presented here. Most of the time, the check is a ``post-hoc'' one, after
obtaining the fitted model.

VIF less than 2 is ideal, and less than 5 is the bottom line.

\texttt{\{r\ comment=\textquotesingle{}\textquotesingle{}\}\ fit\ =\ lm(salary\ \textasciitilde{}\ educ\ +\ salbegin\ +\ jobtime\ +\ prevexp,data\ =\ df)\ car::vif(fit)}

\hypertarget{mlr}{%
\section{MLR}\label{mlr}}

\hypertarget{basic-summary}{%
\subsection{basic summary}\label{basic-summary}}

Firstly, use \texttt{lm()} to get the linear model.

Subsequently, use \texttt{summary()} to get the detailed result.

Period.
\texttt{\{r\ comment=\textquotesingle{}\textquotesingle{}\}\ fit\ =\ lm(salary\ \textasciitilde{}\ educ\ +\ salbegin\ +\ jobtime\ +\ prevexp,data\ =\ df)\ summary(fit)}

It's noteworthy that \texttt{bruceR::print\_table(model)} can spring up
the regression result as a decent table, and can be exported to a
\texttt{.doc} file.

However, fine with \texttt{summary()}, decency doesn't matter much here,
and the \texttt{summary()} can give more information about the model
itself that is essential, such F value and multiple R squared.

\hypertarget{coefficients}{%
\subsection{coefficients}\label{coefficients}}

Use \texttt{coefficients(model)} from base R to see the coefficients for
each IV.

Furthermore, use \texttt{confint(model,level\ =\ 0.95)} from base R to
see the CIs (95\% by default) for each coefficients.

\texttt{\{r\}\ coefficients(fit)\ confint(fit,\ level=0.95)}

Multiple R squared is reported through \texttt{summary()}.

By definition, Multiple R squared is defined as the square of the
correlation between estimated (fitted) value and real values.
\texttt{\{r\ comment=\textquotesingle{}\textquotesingle{}\}\ cor(salary,fit\$fitted.values)}

\hypertarget{model-comparison}{%
\subsection{model comparison}\label{model-comparison}}

Compare the model with the ``zero'' model, formally called the null
model. The null model has no IV but with intercept, which is the mean.
The null model uses the mean to predict the data or explain the
variation of values.

Get the null model by formula \texttt{DV\textasciitilde{}1}.

\texttt{\{r\ comment=\textquotesingle{}\textquotesingle{}\}\ fit2\ =\ lm(salary\ \textasciitilde{}\ 1\ )}

Then use \texttt{anova(model1,model2)} to compare the two models.

\texttt{\{r\ comment=\textquotesingle{}\textquotesingle{}\}\ anova(fit,fit2)}

If we think further, the F test for the MLR model is to test with null
hypothesis that all the coefficients of all the IVs are zero. Therefore,
the model comparison by essence coincides with the F test! Obviously,
the F value ``copies'' that in the \texttt{summary()} result.

\hypertarget{standard-mlr}{%
\section{Standard MLR}\label{standard-mlr}}

Transform the formula with each variable embedded into the function
\texttt{scale()}, no other adjustments!

Under standard MLR, the coefficients are standardized ones.
\texttt{\{r\ comment=\textquotesingle{}\textquotesingle{}\}\ fit\_std\ =\ lm(scale(salary)\ \textasciitilde{}\ scale(educ)\ +\ scale(salbegin)\ +\ scale(jobtime)\ +\ scale(prevexp),data\ =\ df)}

\texttt{\{r\ comment=\textquotesingle{}\textquotesingle{}\}\ summary(fit\_std)}

\hypertarget{semi--partial-for-x}{%
\subsection{(semi-) partial for x}\label{semi--partial-for-x}}

(Semi-) partial correlation indicates the variance accounted for by the
certain IV specifically, which is especially intriguing and important in
standard MlR.

Use \texttt{pcor()} from the package \texttt{ppcor} to compute the
partial correlation for IVs.

Use \texttt{spcor()} from the package \texttt{ppcor} to compute the
semi-partial correlation for IVs.

Throw the regressed variables (both IVs \& DV) into the functions above,
and period.

The results will include estimated value of correlation and its p-value.
However, the drawback is that the returned matrix does not include
rownames and colnames, making it reading-unfriendly. However, we only
care about the (semi-) partial correlation between DV and all IVs, which
means only the very one row or column is enough.

Therefore, DV goes first! Only check the first row or column! ```\{r
comment='\,'\} \# partial correlation df \textbar\textgreater{}
select(salary,educ,salbegin,jobtime,prevexp) \textbar\textgreater{}
ppcor::pcor()

\hypertarget{part-or-semi-partial-correlation}{%
\section{part or semi-partial
correlation}\label{part-or-semi-partial-correlation}}

df \textbar\textgreater{} select(salary,educ,salbegin,jobtime,prevexp)
\textbar\textgreater{} ppcor::spcor()

\begin{verbatim}

# Goodness of Fit
## fitted values & real values
```{r comment=''}
fitted_y = fitted(fit) # predicted values
plot(df$salary,fitted_y,'p')
\end{verbatim}

\hypertarget{residual}{%
\subsection{residual}\label{residual}}

If the model fits well, then the residual should follow a normal
distribution.

Fitted model's residuals can be obtained through
\texttt{residuals(model)} or \texttt{model\$residuals}. This works the
same for other attributes of a linear model, such as coefficients.
\texttt{\{r\ comment=\textquotesingle{}\textquotesingle{}\}\ hist(residuals(fit))\ \#\ hist(fit\$residuals)}

\hypertarget{plot-more}{%
\subsection{plot more}\label{plot-more}}

Use \texttt{plot(model)} to get multiple plots reflecting the goodness
of fit from many perspectives.

Specify parameter \texttt{which\ =\ vector\_of\_range} to designate the
plots to show, 6 the most. \texttt{which\ =\ 1:4} by default.

Following the order, the six plots are:

\begin{itemize}
\tightlist
\item
  Residuals \& Fitted Values

  \begin{itemize}
  \tightlist
  \item
    test the homoscedascitiy assumption
  \end{itemize}
\item
  QQ Plot for Residuals

  \begin{itemize}
  \tightlist
  \item
    test the Gaussian assumption
  \end{itemize}
\item
  Square Root of Standardized Residuals \& Fitted Value

  \begin{itemize}
  \tightlist
  \item
    similar to the first one
  \end{itemize}
\item
  Cook's Distance

  \begin{itemize}
  \tightlist
  \item
    check outliers
  \end{itemize}
\item
  Standardized Residuals \& Leverage

  \begin{itemize}
  \tightlist
  \item
    See if influential points have huge residual, if so, possibly
    outlier!
  \end{itemize}
\item
  Cook's Distance \& Leverage

  \begin{itemize}
  \tightlist
  \item
    Large in both, possibly outlier!
    \texttt{\{r\ comment=\textquotesingle{}\textquotesingle{}\}\ plot(fit,\ which=1:6)}
  \end{itemize}
\end{itemize}

There's another way to do the model checking. Use
\texttt{performance::check\_model(model)}.

The default option is to check all the ``vif'', ``qq'', ``normality'',
``linearity'', ``ncv'', ``homogeneity'', ``outliers'', ``reqq'',
``pp\_check'', ``binned\_residuals'' and ``overdispersion'', which can
be specified through \texttt{check\ =\ specifying\_vector}. Here I'll
take the check of outliers as an example.
\texttt{\{r\ comment=\textquotesingle{}\textquotesingle{}\}\ performance::check\_model(fit,check\ =\ \textquotesingle{}outliers\textquotesingle{})}

Do not forget to \texttt{detach()} the attached data.
\texttt{\{r\ comment=\textquotesingle{}\textquotesingle{}\}\ detach(df)}

\hypertarget{appendix}{%
\section{Appendix}\label{appendix}}

\hypertarget{attach-and-detach}{%
\subsection{Attach and Detach}\label{attach-and-detach}}

We can use the function \texttt{attach()} to use variables in the
database conveniently. After \texttt{attach()}, column variables can be
used directly without \texttt{\$}. (R will do the searching work behind
for you!)

However, sometimes names of the variables will conflict with each other
and cause problems. Not recommended from my perspective, because you'll
not know what database the variable is from, especially when you're
trying to check the codes. Combined with tube-friendly and
formula-pattern genre, this won't be any more trouble.

\texttt{\{r\}\ attach(df)}

After attaching a database, remember to \texttt{detach()} that finally,
in case of conflicts with other variables.

\hypertarget{mlr-1}{%
\subsection{MLR}\label{mlr-1}}

\texttt{\{r\}\ heart\ =\ read.csv(\textquotesingle{}heart.csv\textquotesingle{})}

\hypertarget{fit-the-model}{%
\subsubsection{Fit the Model}\label{fit-the-model}}

\texttt{\{r\}\ mfit\ =\ lm(scale(heart.disease)\textasciitilde{}scale(biking)+scale(smoking),\ \ \ \ \ \ \ \ \ \ \ data\ =\ heart)}

\hypertarget{summary-the-model}{%
\subsubsection{Summary the Model}\label{summary-the-model}}

\texttt{\{r\}\ summary(mfit)}

It's predicted that if \texttt{biking} changes for one standard
deviation, then \texttt{heart.disease} is to change for -0.940 standard
deviation, and that if \texttt{smoking} changes for one standard
deviation, then \texttt{heart.disease} is to change for 0.323 standard
deviation. Meanwhile, coefficients of both \texttt{biking} and
\texttt{smoking} are highly significant with \(p<.001\).

\hypertarget{goodness-of-fit}{%
\subsubsection{Goodness of Fit}\label{goodness-of-fit}}

After qualitative analysis, now move on to quantitative works.

\texttt{\{r\ comment=\textquotesingle{}\textquotesingle{}\}\ result\ =\ summary(mfit)\ result\$adj.r.squared}

\(Adjusted\ R^2\) of 0.9795 indicates a highly awesome goodness of fit
and its ability to explain the variation of the data!

\texttt{\{r\ comment=\textquotesingle{}\textquotesingle{}\}\ cor(x\ =\ fitted(mfit),\ y\ =\ heart\$heart.disease)}

The correlation between the fitted values and the observed (real) values
is as high as 0.9897563, which is exactly the square root of \(R^2\).
Jointly speaking, the model has great credibility in prediction.

\hypertarget{vif}{%
\subsubsection{VIF}\label{vif}}

\texttt{\{r\ comment=\textquotesingle{}\textquotesingle{}\}\ car::vif(mfit)}

VIFs of \texttt{biking} and \texttt{smoking} are 1.000229 and 1.000229
respectively, far less than 5, showing that there should be no worries
about multicollinearity.

\hypertarget{plots-of-residuals}{%
\subsubsection{Plots of Residuals}\label{plots-of-residuals}}

\texttt{\{r\}\ par(mfrow\ =\ c(1,2))\ plot(mfit,which\ =\ c(1,3))}

From the first plot, it's shown that there's no obvious tendency between
the fitted values and the residuals.

\texttt{\{r\}\ par(mfrow\ =\ c(1,2))\ plot(mfit,which\ =\ 2)\ hist(mfit\$residuals)}

In QQ plot, the scatter points of residuals closely center around the
line offered by normal distribution. In the histogram of residuals, a
clear prototype of normal distribution is depicted. What's more, move on
from qualitative analysis to the quantitative.

\texttt{\{r\ comment=\textquotesingle{}\textquotesingle{}\}\ shapiro.test(mfit\$residuals)\ mfit\$residuals\ \textbar{}\textgreater{}\ mean()}

From the Shapiro test, \(W = 0.9969963,p=0.4935143\), indicating that
the residuals from the fitted model follow a normal distribution.
Moreover, the mean is computed to be (extremely close to) zero.

Therefore, we can safely say that the residuals of the fitted model do
follow a normal distribution.

Based on what we've discussed, we can conclude confidently that the
model fits the data well!

\end{document}
